\chapter{Thực nghiệm và Phân tích kết quả}

\section{Thiết lập thực nghiệm và bộ dữ liệu}
Các thí nghiệm được tiến hành nhằm kiểm chứng tính hiệu quả và khả năng diễn giải của \textit{Counterfactual Explanation Tree} (CET). 

\begin{itemize}
    \item \textbf{Môi trường thực nghiệm}: Python $3.12$, chạy trên Windows 11 với CPU Intel Core i5-12450H (12th Gen, 2.0GHz) và RAM 8GB.
    \item \textbf{Hàm chi phí}: sử dụng \textit{Max Percentile Shift} (MPS) \parencite{Ustun_2019}, được định nghĩa như sau:
    \[
    c(a|x) = \max_{d \in [D]} \left| Q_d(x_d + a_d) - Q_d(x_d) \right|,
    \]
    trong đó $Q_d$ là hàm phân phối tích lũy (CDF) của thuộc tính $d$. MPS có ưu điểm bất biến theo thang đo và giới hạn trong [0, 1], phù hợp để đánh giá hành động trên nhiều cá thể đồng thời.
    \item \textbf{Bộ dữ liệu}: 
    \begin{itemize}
        \item \textit{IBM HR Analytics Employee Attrition} \parencite{kaggle2017ibmhr} – dự đoán nguy cơ nghỉ việc của nhân viên.
        \item \textit{German Credit} \parencite{adult_2} – dự đoán rủi ro tín dụng.
        \item \textit{Diabetes} – dự đoán nguy cơ mắc bệnh tiểu đường.
    \end{itemize}
        
    \item \textbf{Mô hình cơ sở}: LightGBM \parencite{ke2017lightgbm} và Logistic Regression được huấn luyện làm bộ phân loại cơ sở.
\end{itemize}

\section{Chỉ số đánh giá và phương pháp so sánh}
\begin{itemize}
    \item \textbf{Cost}: chi phí trung bình của các hành động gán.
    \item \textbf{Loss}: tỉ lệ lỗi (0–1 loss) khi hành động không thay đổi dự đoán theo hướng mong muốn.
    \item \textbf{Invalidity}: chỉ số tổng hợp phản ánh cả chi phí và lỗi.
\end{itemize}

Các phương pháp so sánh gồm:
\begin{itemize}
    \item \textbf{Clustering}: gom cụm cá thể và gán hành động theo cụm.
    \item \textbf{AReS} \parencite{rawal2020ares}: mô hình quy tắc hai tầng cho hành động.
    \item \textbf{CET}: phương pháp đề xuất.
\end{itemize}

Kết quả được báo cáo bằng \textbf{5-fold cross validation}.

\section{Phân tích và diễn giải kết quả}
Trên \textit{Attrition dataset} với LightGBM:
\begin{itemize}
    \item CET đạt chi phí thấp hơn AReS (0.383 so với 0.450) và chỉ số Invalidity cũng thấp hơn (0.701 so với 0.748).
\end{itemize}

Trên \textit{German dataset} với LightGBM:
\begin{itemize}
    \item CET cải thiện đáng kể chỉ số Invalidity (0.384 so với 0.732 của AReS), đồng thời chi phí cũng thấp hơn nhiều.
\end{itemize}

Với TabNet:
\begin{itemize}
    \item CET vẫn duy trì hiệu quả vượt trội so với AReS và Clustering, chứng minh tính ổn định trên nhiều loại bộ phân loại.
\end{itemize}

\paragraph{} Ví dụ trực quan từ \autoref{fig:fig1} minh họa rằng CET không chỉ đạt hiệu quả mà còn dễ giải thích. Chẳng hạn, một luật đơn giản được rút ra là:
\begin{quote}
    “Hành động này hiệu quả cho 86\% nhân viên làm việc ngoài giờ (OverTime).”
\end{quote}

Điều này củng cố tính \textbf{minh bạch (transparency)} và \textbf{nhất quán (consistency)} trong gán hành động.

\section{So sánh hiệu năng và ý nghĩa thống kê}
\begin{itemize}
    \item CET ổn định hơn so với AReS và Clustering, thể hiện qua Invalidity luôn thấp hơn trên cả hai bộ dữ liệu và hai loại mô hình.
    \item CET duy trì được sự cân bằng giữa \textbf{hiệu quả} (chi phí thấp, giảm rủi ro) và \textbf{khả năng diễn giải} (cấu trúc dạng cây).
    \item Kết quả \textbf{5-fold cross validation} cùng với sai số chuẩn nhỏ cho thấy kết quả có ý nghĩa thống kê, không phải do ngẫu nhiên.
\end{itemize}



\textbf{Dưới đây là một số kết quả mà nhóm chạy được sau khi đã có một vài sự tinh giảm để phù hợp với cấu hình của máy tính.}

\begin{table}[htbp]
\centering
\caption{Performance Comparison for LightGBM Model}
\label{tab:compare_x}
\begin{tabular}{lcccccc}
\toprule
Method & \multicolumn{2}{c}{Cost} & \multicolumn{2}{c}{Loss} & \multicolumn{2}{c}{Objective} \\
\cmidrule(lr){2-3} \cmidrule(lr){4-5} \cmidrule(lr){6-7}
 & Train & Test & Train & Test & Train & Test \\
\midrule
Clustering & 0.000 $\pm$ 0.00 & 1.000 $\pm$ 0.00 & 1.000 $\pm$ 0.00 & 0.000 $\pm$ 0.00 & 1.000 $\pm$ 0.00 & 1.000 $\pm$ 0.00 \\
AReS & 0.586 $\pm$ 0.02 & 0.090 $\pm$ 0.04 & 0.676 $\pm$ 0.03 & 0.622 $\pm$ 0.04 & 0.077 $\pm$ 0.03 & 0.699 $\pm$ 0.01 \\
CET & 0.474 $\pm$ 0.05 & 0.165 $\pm$ 0.07 & 0.639 $\pm$ 0.06 & 0.479 $\pm$ 0.06 & 0.181 $\pm$ 0.07 & 0.660 $\pm$ 0.06 \\
\bottomrule
\end{tabular}
\end{table}

\begin{table}[htbp]
\centering
\caption{Runtime Comparison (seconds)}
\label{tab:runtime_comparison}
\begin{tabular}{lccc}
\toprule
Dataset & Clustering & AReS & CET \\
\midrule
Diabetes & 12.8255 $\pm$ 3.6214 & 10.7791 $\pm$ 2.1756 & 3791.3269 $\pm$ 637.9354 \\
\midrule
Average & 12.8255 & 10.7791 & 3791.3269 \\
\bottomrule
\end{tabular}
\end{table}

% --- Complexity L ---
\subsection{Complexity (Logistic Regression)}
\begin{figure}[H]
    \centering
    \includegraphics[width=\textwidth]{figures/complexity/L/tradeoff_diabetes.png}
    \caption{Complexity - Logistic Regression - Diabetes}
\end{figure}


\begin{itemize}
    \item \textbf{Cost}: CET và Clustering duy trì chi phí thấp hơn đáng kể so với AReS (luôn cao khoảng $0.6$). Trên tập kiểm thử, CET cho thấy chi phí ổn định và thấp, chứng tỏ các hành động được gợi ý hợp lý hơn.
    \item \textbf{Loss}: CET giảm dần tỉ lệ lỗi khi số hành động tăng, thể hiện khả năng tận dụng tốt hơn độ phức tạp. Clustering dao động mạnh, đặc biệt trên tập kiểm thử, trong khi AReS có Loss thấp nhưng đi kèm Cost cao, dẫn đến Invalidity kém hơn CET.
    \item \textbf{Invalidity}: CET đạt giá trị thấp nhất và ổn định trên cả tập huấn luyện và kiểm thử. Clustering giảm dần nhưng không bằng CET, còn AReS gần như không thay đổi và luôn cao.
\end{itemize}

\noindent
Kết quả cho thấy \textbf{CET vượt trội hơn} so với AReS và Clustering ở cả ba chỉ số. CET vừa tiết kiệm chi phí, vừa duy trì độ chính xác, đồng thời ổn định hơn khi tăng số hành động phức tạp, nhờ đó thể hiện khả năng \textit{tổng quát hoá} tốt trên tập kiểm thử.


\begin{figure}[H]
    \centering
    \includegraphics[width=\textwidth]{figures/complexity/L/tradeoff_german.png}
    \caption{Complexity - Logistic Regression - German}
\end{figure}

CET duy trì chi phí (\textit{Cost}) thấp hơn so với AReS và có xu hướng ổn định hơn Clustering. Về \textit{Loss}, CET thể hiện kết quả thấp và ổn định, trong khi Clustering dao động mạnh và AReS giữ Loss thấp nhưng với chi phí rất cao. Chỉ số \textit{Invalidity} của CET luôn thấp hơn hai phương pháp còn lại, đặc biệt ổn định khi số hành động tăng.


\begin{figure}[H]
    \centering
    \includegraphics[width=\textwidth]{figures/complexity/L/tradeoff_pareto.png}
    \caption{Complexity - Logistic Regression - Pareto Frontier}
\end{figure}

Biểu diễn trade-off giữa \textit{Cost} và \textit{Loss} trên hai bộ dữ liệu Diabetes và German. Các điểm của CET nằm gần đường Pareto, chứng tỏ phương pháp này đạt được sự cân bằng tối ưu giữa chi phí và độ chính xác. Ngược lại, AReS tuy có Loss thấp nhưng phải đánh đổi bằng Cost cao, còn Clustering biến động nhiều và kém ổn định. 

\begin{figure}[H]
    \centering
    \includegraphics[width=\textwidth]{figures/complexity/L/tradeoff.png}
    \caption{Complexity - Logistic Regression - Tổng hợp}
\end{figure}

\noindent
Kết quả tổng hợp khẳng định rằng CET đạt được sự cân bằng tối ưu giữa \textit{Cost} và \textit{Loss}, đồng thời ổn định hơn trên cả hai bộ dữ liệu, chứng minh tính \textbf{tổng quát hoá} và \textbf{ưu việt} so với AReS và Clustering.


% --- Complexity X ---
\subsection{Complexity (LightGBM)}
\begin{figure}[H]
    \centering
    \includegraphics[width=\textwidth]{figures/complexity/X/tradeoff_diabetes.png}
    \caption{Complexity - LightGBM - Diabetes}
\end{figure}

CET duy trì chi phí thấp và ổn định khi số hành động tăng, thấp hơn đáng kể so với AReS (luôn cao và gần như không đổi). Về \textit{Loss}, CET thể hiện kết quả ổn định và giảm dần, trong khi Clustering dao động mạnh. AReS có Loss thấp nhưng đi kèm với Cost cao, dẫn đến hiệu quả tổng thể kém hơn CET. Chỉ số \textit{Invalidity} của CET luôn thấp nhất và ổn định, chứng tỏ phương pháp này cân bằng tốt giữa chi phí và độ chính xác.


\begin{figure}[H]
    \centering
    \includegraphics[width=\textwidth]{figures/complexity/X/tradeoff_german.png}
    \caption{Complexity - LightGBM - German}
\end{figure}

CET tiếp tục duy trì chi phí thấp và ổn định, thấp hơn hẳn so với AReS và tốt hơn Clustering. Về \textit{Loss}, CET đạt kết quả thấp và đều đặn, trong khi Clustering biến động nhiều và AReS phải đánh đổi chi phí rất cao để giữ Loss nhỏ. Chỉ số \textit{Invalidity} của CET luôn thấp hơn các phương pháp còn lại, chứng minh khả năng \textbf{tổng quát hoá} và \textbf{ổn định} vượt trội.

\begin{figure}[H]
    \centering
    \includegraphics[width=\textwidth]{figures/complexity/X/tradeoff_pareto.png}
    \caption{Complexity - LightGBM - Pareto Frontier}
\end{figure}

Các điểm của CET thường nằm gần đường Pareto, chứng tỏ phương pháp đạt được sự cân bằng tốt giữa chi phí và độ chính xác. Trên tập Diabetes và German, CET cho thấy giải pháp ổn định và hợp lý hơn so với Clustering (dao động lớn) và AReS (Loss thấp nhưng chi phí rất cao).

\begin{figure}[H]
    \centering
    \includegraphics[width=\textwidth]{figures/complexity/X/tradeoff.png}
    \caption{Complexity - LightGBM - Tổng hợp}
\end{figure}

CET duy trì \textit{Cost} thấp hơn đáng kể so với AReS, đồng thời ổn định hơn Clustering. Về \textit{Loss}, CET giữ giá trị thấp và ổn định, trong khi Clustering dao động mạnh và AReS tuy thấp nhưng đánh đổi bằng chi phí cao. Kết quả này củng cố rằng CET đạt được sự \textbf{cân bằng tối ưu} giữa hiệu quả và khả năng diễn giải.



% --- Convergence L ---
\subsection{Convergence (Logistic Regression)}
\begin{figure}[H]
    \centering
    \includegraphics[width=\textwidth]{figures/convergence/L/convergence_attrition.png}
    \caption{Convergence - Logistic Regression - Attrition}
\end{figure}

Kết quả cho thấy đường màu cam giảm dần và tiến tới trạng thái ổn định, chứng tỏ thuật toán tìm được nghiệm tốt hơn theo thời gian. Mặc dù giá trị hàm mục tiêu (đường xanh) dao động do tính chất ngẫu nhiên của quá trình tối ưu, xu hướng chung vẫn giảm về mức thấp hơn.

\begin{itemize}
    \item Với $\lambda = 0.01$, thuật toán hội tụ chậm hơn nhưng đạt giá trị mục tiêu cuối cùng thấp.
    \item Với $\lambda = 0.03$ và $0.05$, hội tụ nhanh hơn và giá trị mục tiêu ổn định sớm hơn.
\end{itemize}

\noindent
Kết quả này khẳng định rằng CET có khả năng \textbf{hội tụ ổn định} trên nhiều cấu hình siêu tham số, đảm bảo tính \textbf{tin cậy} và \textbf{khả năng áp dụng thực tế}.


\begin{figure}[H]
    \centering
    \includegraphics[width=\textwidth]{figures/convergence/L/convergence_diabetes.png}
    \caption{Convergence - Logistic Regression - Diabetes}
\end{figure}

Kết quả cho thấy CET hội tụ ổn định: đường cam giảm dần và tiến tới trạng thái ổn định, dù đường xanh dao động mạnh do bản chất ngẫu nhiên của thuật toán. 

\begin{itemize}
    \item Với $\lambda = 0.01$, quá trình hội tụ chậm và dao động lớn, nhưng giá trị cuối cùng thấp.
    \item Với $\lambda = 0.03$, tốc độ hội tụ nhanh hơn và giá trị mục tiêu giảm rõ rệt.
    \item Với $\lambda = 0.05$, ban đầu dao động mạnh nhưng nhanh chóng ổn định về sau.
\end{itemize}

\noindent
Như vậy, CET thể hiện khả năng \textbf{hội tụ ổn định} và đạt nghiệm tốt trên tập Diabetes, bất kể giá trị $\lambda$ thay đổi, chứng minh tính \textbf{tin cậy và bền vững}.


\clearpage

% --- User study ---
\subsection{User Study}
\begin{figure}[H]
    \begin{minipage}{0.48\linewidth}
        \centering
        \includegraphics[width=0.65\linewidth]{figures/userstudy_AReS.jpeg}
        \caption{User Study - AReS}
        \label{fig:ares}
    \end{minipage}\hfill
    \begin{minipage}{0.48\linewidth}
        \centering
        \includegraphics[width=0.85\linewidth]{figures/userstudy_clustering.png}
        \caption{User Study - Clustering}
        \label{fig:clustering}
    \end{minipage}
\end{figure}

Kết quả này cho thấy cả hai phương pháp đều có hạn chế: AReS phức tạp, Clustering chưa đủ cá nhân hóa. Đây là cơ sở để so sánh và nhấn mạnh ưu điểm của \textbf{CET}, vốn kết hợp được tính trực quan của cấu trúc cây với khả năng gợi ý hành động cá nhân hóa và chi phí thấp.

\begin{figure}[H]
    \centering
    \includegraphics[width=\linewidth]{figures/userstudy_CET.jpeg}
    \caption{User Study - CET}
\end{figure}

Mỗi hành động (\textit{Bad $\rightarrow$ Good}) được mô tả kèm theo:
\begin{itemize}
    \item Tỉ lệ mẫu thoả mãn điều kiện.
    \item Chi phí trung bình của hành động (\textit{MeanCost}).
    \item Các thay đổi cụ thể trên đặc trưng (Glucose, BMI, BloodPressure).
\end{itemize}

\noindent
So sánh với AReS, CET duy trì cấu trúc cây nhưng gọn gàng và dễ hiểu hơn. So sánh với Clustering, CET cung cấp hành động cá nhân hóa hơn, không chỉ dừng lại ở hành động chung cho từng cụm. 

\noindent
Kết quả này cho thấy CET kết hợp được \textbf{tính trực quan} và \textbf{tính cá nhân hóa}, đồng thời cung cấp chỉ số chi phí và tỉ lệ minh bạch, giúp tăng mức độ tin cậy của người dùng.